\begin{question}
\section*{Theory}

\begin{enumerate}
\item State the necessary assumptions that are needed to study Rossby waves. 

\item Write down the definition of the barotropic stream function. 

\item Write down the equation for the stream function when assuming boundary condition $\Psi_s(x_e)=0$. State the name of the circulation that is described by the equation. 

\item Describe the shape of the gyre. 

\item State the additional assumptions for Stommel’s theory. 

\item Describe the main changes in the shape of the stream function. 

\item Give the equation for Stommel’s solution. 

\item Describe the shape of the wind stress (discuss the equation for the wind stress) and the geographical situation where this can commonly be applied. 
\end{enumerate}

\section*{Experiments}
\subsection*{Sverdrup circulation}

\begin{enumerate}
\item Provide plots of the stream function for all your experiments at appropriate time steps.

\item What effect does $\beta = \partial f / \partial y$ have on the system? 

\item What effect does the wind forcing have? 

\item Provide cross sections (in $x$-direction) of the stream functions computed from your experiments. 

\item From all your plots, discuss if the results resemble theory. Why / why not? 

\item Do the results improve if the size of the domain is increased? 

\item What are the limitations of the Sverdrup model? 
\end{enumerate}

\subsection*{Stommel’s theory}

\begin{enumerate}
\item Provide plots of the stream functions of your experiments. 

\item Describe how the results have changed from Sverdrup circulation to Stommel’s theory. 

\item Do theory and numerical solution correlate with each other? 

\item Provide:
\begin{itemize}
\item cross sections (in $x$-direction) of the stream functions computed from your experiments. 
\item the difference of the cross section of stream functions between theory and simulation results. 
\end{itemize}
and make sure that they also show the Stommel thickness from theory. 

\item From that, discuss the role of the friction coefficient for the shape. 
\end{enumerate}
\end{question}